\section{Using the original ablations to design broader tests}
\label{sec:cross_method_design}
This section specifies prospective comparisons, not completed results. The method list is grounded in the full TextSeal2 and published DistSeal papers~\citep{textseal2,distseal}. Additional methods enter only after the primary defense demonstrates a benefit against a preservation-qualified attack. Every additional primary paper must be read in full, and its official code, checkpoint, license and supported interface checked before implementation. A locally available wrapper alone does not establish a faithful baseline.

\paragraph{Text comparisons.} TextSeal, SynthID-Text and textGrain are the first three candidates because TextSeal2 compares them directly. KGW supplies a fourth, classical family if the final interface defense supports it. The TextSeal2 textGrain implementation is explicitly a reimplementation, with choices that may differ from unreleased original code. AutoMark is retained as a diagnostic: without speculative decoding, its copula and TextSeal2's virtual-key construction agree in distribution. Counting these two as independent families would overstate coverage. Native speculative decoding, if tested, forms a separate experiment.

Table~6 of TextSeal2 also motivates a detector ablation: compare the native key-only test with entropy weighting, calibrating each on a separate null set. Its Tables~7--8 motivate checking watermark strength and context length. In particular, a ladder of window lengths is already studied and can lose robustness under edits; it is neither an untested universal improvement nor a new contribution here. Any additional context design must therefore be compared with the existing short-window and multi-window controls at matched diversity. These controls are frozen before final attack outputs are inspected.

\paragraph{Image comparisons.} DistSeal Appendix~E provides the most directly compatible route for decoder experiments: WAM, TrustMark, CIN and MBRS can serve as four distinct post-hoc teachers distilled into the same DC-AE decoder. The released implementation and weights must first be verified. This would test transfer of a hardening procedure across watermark teachers while holding the decoder architecture and attacker access fixed. It would not establish that the original post-hoc method, or an untouched external marking layer, withstands a weight attack.

The initial-quality gate is essential. DistSeal's RAR decoder results include failed distillations for several otherwise valid pixel watermarks, so those cases cannot be presented as successful watermark removal. Stable Signature is another candidate, but DistSeal Appendix~F.1 applies a Stable Signature-style loss to the same latent teacher. That loss ablation is distinct from evaluating an independently released Stable Signature model. RoSteALS and Tree-Ring supply additional comparisons only under explicitly compatible access settings; a zero-bit initial-noise watermark cannot silently be treated as a learned in-weight payload.

\paragraph{Ablations that identify the cause of an improvement.} DistSeal's teacher and extractor-loss comparisons motivate separating the choice of marked target from the loss used to retain quality. For the present image candidate, a crossed experiment should vary fixed versus adapted targets and denoising versus prediction retention at matched coefficients. The completed development revision changes both retention type and coefficient; it cannot identify their separate effects. Useful subsequent controls include an adapted clean reference without marking, a marked target without inner adaptation, and independently drawn inner tasks. Generic teacher guidance, watermark composition and meta-learning are already prior art.

\paragraph{Shared evaluation across modalities.} The common question is whether a defense keeps its signal after an allowed change while preserving useful content. Within each modality, all methods use the same content identities, reserved splits, attacker information, update or query budget, and quality-only selection rule. Report baseline detection and payload length alongside paired quality, raw removal and joint removal; threshold calibration uses a separate null set at a common nominal false-positive rate. Repeated keys on the same null outputs do not create independent null texts. Fresh final prompts, seeds and classes are reserved before tuning, and paired uncertainty is reported. A failure to find any eligible quality-preserving attack is reported explicitly, without being converted into evidence of a successful robust defense. All added runs and confirmations share the 60 H100-hour cap.
